\documentclass[11pt]{article}
\usepackage[a4paper,margin=1in]{geometry}
\usepackage[T1]{fontenc}
\usepackage{lmodern}
\usepackage{amsmath,amssymb,amsthm,booktabs,enumitem,mathtools,microtype}
\usepackage[numbers,sort&compress]{natbib}
\usepackage[hidelinks]{hyperref}

\newtheorem{theorem}{Theorem}[section]
\newtheorem{proposition}[theorem]{Proposition}
\newtheorem{lemma}[theorem]{Lemma}
\newtheorem{corollary}[theorem]{Corollary}
\theoremstyle{definition}
\newtheorem{definition}[theorem]{Definition}
\theoremstyle{remark}
\newtheorem{remark}[theorem]{Remark}

\newcommand{\mob}{\mu}
\newcommand{\Z}{\mathbb Z}
\newcommand{\cA}{\mathcal A}
\newcommand{\cB}{\mathcal B}
\newcommand{\cC}{\mathcal C}
\newcommand{\DeltaY}{\Delta_y}

\hypersetup{
  pdftitle={Prime-Tail Scale Separation below the Quadratic Square-Clock Gap},
  pdfauthor={RH research program},
  pdfsubject={Fixed-r prime-square tail asymptotics and normalized square-clock gap limits},
  pdfkeywords={prime number theorem, Abel summation, prime-square tail, scale separation, Euler product}
}

\title{Prime-Tail Scale Separation below\
the Quadratic Square-Clock Gap}
\author{RH research program}
\date{August 8, 2026}

\begin{document}
\maketitle

\begin{abstract}
For the square-clock sequence
$q_y=4\prod_{i\le y}p_i^2$, where $3=p_1<p_2<\cdots$ are the odd
primes, define the strict prime tails
\[
 P_r(y)=\sum_{p>p_y}\frac{1}{(p^2-1)^r},\qquad
 T_y=P_1(y),\qquad S_y=P_2(y).
\]
We prove from the prime number theorem and strict Stieltjes partial
summation that, for every fixed integer $r\ge1$,
\[
 P_r(y)\sim\frac{1}{(2r-1)p_y^{2r-1}\log p_y}.
\]
Consequently, for each fixed partition
$\lambda=1^{k_1}\cdots d^{k_d}$ of degree $d$ and length $\ell$,
\[
 \prod_rP_r(y)^{k_r}
 \sim\prod_r(2r-1)^{-k_r}
 p_y^{-(2d-\ell)}(\log p_y)^{-\ell}.
\]
In particular $T_y^3=o(S_y)$ and $S_y=o(T_y^2)$, with
$S_y/[T_y^3(\log p_y)^2]\to1/3$.

Combining this scale dictionary with the exact RH-382 input
\[
 B_\infty-G(q_y)=A T_y+B T_y^2+C S_y+O(T_y^3)
\]
gives five normalized gap limits, including two equivalent normalizations
of the independent $S_y$ coefficient and a logarithmically amplified
third-order residual.  The subtractions remain the exact intrinsic terms
$A T_y$ and $B T_y^2$: bare prime-number-theorem equivalents cannot be
substituted at the smaller scales.  A precision-$80$ directed-rounding
certificate proves
\[
 1.5463476716710499204
 \le Y_\infty-2m_\infty
 \le1.5484488989771761113,
\]
so $C=(Y_\infty-2m_\infty)/\pi^2>0$.  Hence the twice-subtracted residual
is eventually positive and its quotient by $T_y^3$ tends to $+\infty$,
without an effective threshold.  The result is confined to fixed $r$,
fixed partitions, fixed finite clocks before the prefix limit, and the
universally safe phasewise class with $c_{11}=0$.  It constructs no
operator or zeta-divisor identity and has no implication for the Riemann
hypothesis.
\end{abstract}

\noindent\textbf{Keywords:} prime number theorem; Abel summation;
prime-square Euler tails; fixed partitions; scale separation; finite
clocks.

\section{Frozen class, strict tails, and inherited input}

We work only in the fixed finite-clock class established in the preceding
square-clock papers \citep{RH374,RH379,RH380,RH381,RH382,RH383}.  Fix the
clock $q$ before sending the prefix length $N$ to infinity.  At every phase
$a\in\Z/q\Z$, the sign rule depends on
$(\mob(n-2),\mob(n))\in\{-1,0,1\}^2$, is universally distance-two-safe,
and has vanishing phasewise interpolation coefficient $c_{11}(a)$.  The
resulting optimum is $G(q)$.  An active $c_{11}$ would expose an uncontrolled
phase-weighted shift-two M\"obius correlation and is not admitted.

Let
\[
 3=p_1<p_2<\cdots,\qquad
 q_y=4\prod_{i\le y}p_i^2,\qquad
 a_{j+1}=\frac{1}{p_{j+1}^2-1}.
\]
The intrinsic power-sum coordinates introduced in RH-383 are
\begin{equation}
 P_r(y)=\sum_{j\ge y}a_{j+1}^r
       =\sum_{\substack{p\ \mathrm{prime}\\p>p_y}}
        \frac{1}{(p^2-1)^r},
 \qquad T_y=P_1(y),\qquad S_y=P_2(y).
 \label{eq:tails}
\end{equation}
Thus the endpoint is strict, and the first atom is at $p_{y+1}$.
Absolute convergence makes the exact successor identity
\begin{equation}
 P_r(y)=\frac{1}{(p_{y+1}^2-1)^r}+P_r(y+1)
 \label{eq:successor}
\end{equation}
valid for every $r\ge1$ and $y\ge1$.

Define the limiting Euler ratios
\[
 u_m=\prod_{p\ \mathrm{odd}}\frac{p^2-m}{p^2-1}
 \qquad(2\le m\le8)
\]
and the three inherited linear forms
\begin{align}
 X_\infty&=2u_4-4u_5+6u_6-8u_7+10u_8,\\
 Y_\infty&=6u_4-16u_5+30u_6-48u_7+70u_8,\\
 m_\infty&=2u_3-4u_4+6u_5-8u_6+10u_7-12u_8.
 \label{eq:XYm}
\end{align}
RH-382 proves, within the frozen class and after the fixed-clock prefix
limit has been taken,
\begin{equation}
 \DeltaY:=B_\infty-G(q_y)
 =A T_y+B T_y^2+C S_y+O(T_y^3),
 \label{eq:rh382-input}
\end{equation}
where
\begin{equation}
 A=\frac{2X_\infty}{\pi^2},\qquad
 B=\frac{Y_\infty+2m_\infty}{\pi^2},\qquad
 C=\frac{Y_\infty-2m_\infty}{\pi^2}.
 \label{eq:ABC}
\end{equation}
Equation \eqref{eq:rh382-input} is an immutable predecessor theorem, not a
conclusion fitted from the finite rows in Section~\ref{sec:artifact}.

The prime-counting input is the classical prime number theorem
\begin{equation}
 \pi(x)\sim\frac{x}{\log x},
 \label{eq:pnt}
\end{equation}
for which we use Montgomery and Vaughan, Chapter~6
\citep{MontgomeryVaughan2007}.  The repository provenance for this exact
input is the frozen RH-2 release \citep{RH2}; later M\"obius/Mertens
consequences are not treated as a substitute prime-counting source.

\section{Fixed-r Abel asymptotics}

The boundary term in partial summation is decisive for the leading
constant.  We record it before making any approximation.

\begin{lemma}[Strict Stieltjes identity]
\label{lem:stieltjes}
For $x=p_y$, $r\ge1$, and $f_r(t)=(t^2-1)^{-r}$,
\begin{equation}
 P_r(y)=\int_{(x,\infty)}f_r(t)\,d\pi(t)
 =-\frac{\pi(x)}{(x^2-1)^r}
   -\int_x^\infty \pi(t)f_r'(t)\,dt.
 \label{eq:stieltjes}
\end{equation}
In particular, the negative endpoint term is present because the sum is
over $p>x$.
\end{lemma}

\begin{proof}
Apply Stieltjes integration by parts on $(x,R]$.  Since
$f_r(R)\pi(R)=O(R^{1-2r})\to0$, passage to $R\to\infty$ gives
\eqref{eq:stieltjes}.  The value $\pi(x)$ in the boundary term removes the
atom at $x$ and therefore implements the strict endpoint.
\end{proof}

\begin{theorem}[Fixed prime-power tail]
\label{thm:fixed-r}
For every fixed integer $r\ge1$,
\begin{equation}
 \boxed{
 P_r(y)\sim
 \frac{1}{(2r-1)p_y^{2r-1}\log p_y}}
 \qquad(y\to\infty).
 \label{eq:fixed-r}
\end{equation}
No assertion of uniformity in a growing $r$ is made.
\end{theorem}

\begin{proof}
Put $x=p_y$ and $s=2r$.  For fixed $r$,
\[
 (p^2-1)^{-r}-p^{-2r}=O_r(p^{-2r-2}),
\]
uniformly over primes $p>x$.  Hence
\begin{equation}
 P_r(y)-\sum_{p>x}p^{-s}
 =O_r\!\left(\sum_{n>x}n^{-s-2}\right)
 =O_r(x^{-s-1})
 =o\!\left(\frac{x^{1-s}}{\log x}\right).
 \label{eq:weight-reduction}
\end{equation}
Partial summation, equivalently Lemma~\ref{lem:stieltjes} for $t^{-s}$,
gives the exact strict identity
\begin{equation}
 \sum_{p>x}p^{-s}
 =-\pi(x)x^{-s}+s\int_x^\infty\pi(t)t^{-s-1}\,dt.
 \label{eq:abel-power}
\end{equation}

Set $g(t)=t/\log t$ and
\[
 \varepsilon_x=\sup_{t\ge x}
 \left|\frac{\pi(t)}{g(t)}-1\right|.
\]
The PNT implies $\varepsilon_x\to0$.  This tail supremum is the needed
uniformization: the error in replacing $\pi(t)$ by $g(t)$ in the integral
of \eqref{eq:abel-power} is at most
\[
 s\varepsilon_x\int_x^\infty\frac{t^{-s}}{\log t}\,dt.
\]
For fixed $s>1$, integration by parts (or l'H\^opital's rule) gives
\begin{equation}
 I_s(x):=\int_x^\infty\frac{t^{-s}}{\log t}\,dt
 \sim\frac{x^{1-s}}{(s-1)\log x}.
 \label{eq:Is}
\end{equation}
Substitution in \eqref{eq:abel-power} therefore yields
\[
 \sum_{p>x}p^{-s}
 =\left(-1+\frac{s}{s-1}+o(1)\right)
   \frac{x^{1-s}}{\log x}
 =\left(\frac{1}{s-1}+o(1)\right)
   \frac{x^{1-s}}{\log x}.
\]
Combining this with \eqref{eq:weight-reduction} and $s=2r$ proves
\eqref{eq:fixed-r}.  The negative Abel boundary is precisely what changes
$2r/(2r-1)$ into $1/(2r-1)$.
\end{proof}

\begin{remark}[Endpoint interface versus leading equivalent]
Replacing $p>p_y$ by $p\ge p_y$ adds one atom of order $p_y^{-2r}$,
which is smaller than the main tail by a factor of order
$\log p_y/p_y$.  Likewise $p_{y+1}/p_y\to1$ follows from the PNT.  Thus
the leading equivalent \eqref{eq:fixed-r} cannot reject either an inclusive
endpoint or an RHS written at $p_{y+1}$.  The strict convention and first
atom in \eqref{eq:successor} are frozen by exact membership identities;
they are not advertised as distinct leading asymptotics.
\end{remark}

\section{Fixed partitions and the scale dictionary}

Let $\lambda=1^{k_1}\cdots d^{k_d}$ be a fixed partition, with degree and
length
\[
 |\lambda|=\sum_{r=1}^d rk_r=d,
 \qquad \ell(\lambda)=\sum_{r=1}^d k_r=\ell,
\]
and put $P_\lambda(y)=\prod_rP_r(y)^{k_r}$.

\begin{corollary}[Fixed-partition scale compiler]
\label{cor:partition}
For every fixed partition $\lambda$,
\begin{equation}
 \boxed{
 P_\lambda(y)\sim
 \left(\prod_r(2r-1)^{-k_r}\right)
 p_y^{-(2d-\ell)}(\log p_y)^{-\ell}.}
 \label{eq:partition}
\end{equation}
\end{corollary}

\begin{proof}
Multiply the finitely many fixed-$r$ equivalents in
Theorem~\ref{thm:fixed-r}.  The exponent identity is
$\sum_r(2r-1)k_r=2d-\ell$, while the logarithmic exponent is
$\sum_rk_r=\ell$.
\end{proof}

The corollary is pointwise in $\lambda$.  It does not license a partition
degree or length which grows with $y$.  For $T_y=P_1(y)$ and
$S_y=P_2(y)$ it gives the complete scale ledger needed below.

\begin{proposition}[Five intrinsic scale relations]
\label{prop:scales}
As $y\to\infty$,
\begin{align}
 p_y\log p_y\,T_y&\longrightarrow1,
 \label{eq:scale1}\\
 3p_y^3\log p_y\,S_y&\longrightarrow1,
 \label{eq:scale2}\\
 \frac{S_y}{T_y^2}&\longrightarrow0,
 \label{eq:scale3}\\
 \frac{T_y^3}{S_y}&\longrightarrow0,
 \label{eq:scale4}\\
 \frac{S_y}{T_y^3(\log p_y)^2}&\longrightarrow\frac13.
 \label{eq:scale5}
\end{align}
More explicitly, the last three comparisons have equivalents
\[
 \frac{S_y}{T_y^2}\sim\frac{\log p_y}{3p_y},\qquad
 \frac{T_y^3}{S_y}\sim\frac{3}{(\log p_y)^2},\qquad
 \frac{S_y}{T_y^3(\log p_y)^2}\sim\frac13.
\]
\end{proposition}

\begin{proof}
Theorem~\ref{thm:fixed-r} gives
$T_y\sim[p_y\log p_y]^{-1}$ and
$S_y\sim[3p_y^3\log p_y]^{-1}$.  The displayed relations follow by
division.
\end{proof}

In particular, $S_y$ lies strictly between the two quadratic/cubic
monomials on the asymptotic scale:
\begin{equation}
 T_y^3=o(S_y),\qquad S_y=o(T_y^2).
 \label{eq:scale-order}
\end{equation}
The $P_2$ term in \eqref{eq:rh382-input} is therefore an independent
intermediate scale; it cannot be absorbed into either adjacent monomial.

\section{Five normalized square-clock gap limits}

We now combine only the exact predecessor expansion
\eqref{eq:rh382-input} and Proposition~\ref{prop:scales}.

\begin{theorem}[Normalized gap ledger]
\label{thm:gap-limits}
Within the frozen phasewise class,
\begin{align}
 p_y\log p_y\,\DeltaY&\longrightarrow A,
 \label{eq:L1}\\
 (p_y\log p_y)^2(\DeltaY-A T_y)&\longrightarrow B,
 \label{eq:L2}\\
 \frac{\DeltaY-A T_y-BT_y^2}{S_y}&\longrightarrow C,
 \label{eq:L3}\\
 3p_y^3\log p_y\,(\DeltaY-A T_y-BT_y^2)&\longrightarrow C,
 \label{eq:L4}\\
 \frac{\DeltaY-A T_y-BT_y^2}
 {T_y^3(\log p_y)^2}&\longrightarrow\frac C3.
 \label{eq:L5}
\end{align}
Here $A,B,C$ are exactly those in \eqref{eq:ABC}.
\end{theorem}

\begin{proof}
For \eqref{eq:L1}, multiply \eqref{eq:rh382-input} by
$p_y\log p_y$; the leading term tends to $A$ and all smaller terms vanish.
For \eqref{eq:L2}, exact subtraction of $AT_y$ leaves
$BT_y^2+CS_y+O(T_y^3)$.  Multiplication by
$(p_y\log p_y)^2$ sends the first term to $B$, while
$S_y/T_y^2\to0$ and $T_y\to0$ dispose of the rest.

Put
\begin{equation}
 R_y=\DeltaY-A T_y-BT_y^2.
 \label{eq:residual}
\end{equation}
Then \eqref{eq:rh382-input} says $R_y=CS_y+O(T_y^3)$.
Because $T_y^3/S_y\to0$, division by $S_y$ proves \eqref{eq:L3}.
Equation \eqref{eq:L4} follows from
$3p_y^3\log p_y\,S_y\to1$.  Finally,
\[
 \frac{R_y}{T_y^3(\log p_y)^2}
 =C\frac{S_y}{T_y^3(\log p_y)^2}
   +O\!\left(\frac{1}{(\log p_y)^2}\right),
\]
and \eqref{eq:scale5} proves \eqref{eq:L5}.
\end{proof}

\begin{remark}[Exact-subtraction firewall]
The terms in \eqref{eq:L2}--\eqref{eq:L5} are intrinsic and exact.  Bare
PNT gives only
\[
 T_y=\frac{1+o(1)}{p_y\log p_y},\qquad
 T_y^2=\frac{1+o(1)}{p_y^2(\log p_y)^2}.
\]
It does not say that
$T_y-[p_y\log p_y]^{-1}=o(T_y^2)$, so replacing $AT_y$ by its PNT
surrogate is unlicensed already in \eqref{eq:L2}.  At the still smaller
$S_y$ scale, an $o(T_y^2)$ error can dominate $S_y$; consequently the
exact $BT_y^2$ in \eqref{eq:L3}--\eqref{eq:L5} cannot be replaced by
$B/[p_y^2(\log p_y)^2]$ either.  No effective PNT remainder is silently
assumed.
\end{remark}

\section{A directed interval and the sign of the residual}
\label{sec:interval}

It remains to certify the sign of $C$.  Let $N=100000$, and for
$2\le m\le8$ define the finite product
\[
 U_m(N)=\prod_{\substack{p\le N\\p\ \mathrm{odd}}}
 \frac{p^2-m}{p^2-1}.
\]
For a tail prime put $b_p=(m-1)/(p^2-1)$.  Bonferroni's elementary
product bound gives
\[
 1-\sum_{p>N}b_p\le\prod_{p>N}(1-b_p)\le1.
\]
Moreover
\begin{equation}
 \sum_{p>N}\frac1{p^2-1}
 \le\sum_{n>N}\frac1{n^2-1}
 =\frac12\left(\frac1N+\frac1{N+1}\right)
 =\frac{200001}{20000200000}=:\theta_N.
 \label{eq:theta}
\end{equation}
Therefore the exact enclosure
\begin{equation}
 U_m(N)\bigl(1-(m-1)\theta_N\bigr)
 \le u_m\le U_m(N)
 \label{eq:um-interval}
\end{equation}
holds for every $2\le m\le8$.

The artifact evaluates the $9591$ odd-prime factors at decimal precision
$80$, with floor rounding for lower endpoints, ceiling rounding for upper
endpoints, and binary-float conversion trapped.  In particular, the tail
loss $(m-1)\theta_N$ is first rounded upward under ceiling rounding; only
then is its complement $1-(m-1)\theta_N$ formed under floor rounding for
the lower product endpoint.  The cutoff anchor is the
integer $N=100000$; the last prime below it is $99991$ and the next prime is
$100003$, neither of which replaces $N$ in \eqref{eq:theta}.  Sign-aware
interval arithmetic is applied independently to
\begin{align*}
 Y_\infty&=6u_4-16u_5+30u_6-48u_7+70u_8,\\
 m_\infty&=2u_3-4u_4+6u_5-8u_6+10u_7-12u_8,\\
 Y_\infty-2m_\infty&=-4u_3+14u_4-28u_5+46u_6-68u_7+94u_8.
\end{align*}

\begin{proposition}[Positive contrast certificate]
\label{prop:positive}
The full precision-$80$ endpoints are stored in the exact result ledger.
They imply the shorter outward bounds
\[
 1.5463476716710499204067249
 <Y_\infty-2m_\infty
 <1.5484488989771761112886458.
\]
In particular, outward quantization to $19$ decimal places gives
\begin{equation}
 \boxed{1.5463476716710499204
 \le Y_\infty-2m_\infty
 \le1.5484488989771761113},
 \label{eq:published-interval}
\end{equation}
and $C>0$.
\end{proposition}

\begin{proof}
Equation \eqref{eq:um-interval} supplies rigorous intervals for
$u_2,\ldots,u_8$.  For each coefficient in the last linear form, the lower
endpoint uses the lower $u_m$ endpoint when the coefficient is positive and
the upper endpoint when it is negative; the choices reverse for the upper
endpoint.  Directed arithmetic gives the displayed raw bounds.  Outward
quantization gives \eqref{eq:published-interval}.  Division by $\pi^2>0$
then proves $C>0$.
\end{proof}

\begin{corollary}[Eventual sign and logarithmic amplification]
\label{cor:sign}
The residual $R_y$ in \eqref{eq:residual} is eventually positive and
\begin{equation}
 \frac{R_y}{T_y^3}\longrightarrow+\infty.
 \label{eq:amplification}
\end{equation}
No effective index from which positivity holds is asserted.
\end{corollary}

\begin{proof}
Theorem~\ref{thm:gap-limits} gives $R_y/S_y\to C>0$, proving eventual
positivity.  Also
\[
 \frac{R_y}{T_y^3}=C\frac{S_y}{T_y^3}+O(1)
 =\left(\frac{C}{3}+o(1)\right)(\log p_y)^2+O(1)
 \longrightarrow+\infty.
\]
The PNT input is qualitative, so the proof supplies no effective threshold.
\end{proof}

\section{Executable certificate and immutable provenance}
\label{sec:artifact}

The executable artifact separates symbolic theorem content from finite
reproduction.  The exact grid contains the rows in
Table~\ref{tab:certificate}.  Partition rows enumerate all $66$ partitions
of degrees $1$ through $8$.  Successor rows use
$y\in\{1,2,3,5,8,13\}$ and $r=1,\ldots,8$; they reproduce the exact finite
identity \eqref{eq:successor}, not the $y\to\infty$ theorem.

\begin{table}[ht]
\centering
\small
\begin{tabular}{@{}lrrp{0.50\linewidth}@{}}
\toprule
ledger & rows & arithmetic & role \\
\midrule
fixed $r$ & 8 & rational/symbolic & constants and exponents for $r\le8$ \\
fixed partitions & 66 & rational/symbolic & all partitions through degree $8$ \\
strict successors & 48 & exact fractions & endpoint and first-atom interface \\
scale relations & 5 & algebraic & constants and relative orders \\
normalized gaps & 5 & algebraic & exact-subtraction ledger \\
numeric intervals & 10 & directed decimal & $u_2,\ldots,u_8,Y,m,Y-2m$ \\
negative mutations & 20 & fail-closed & theorem, source, type, JSON, and numeric attacks \\
\bottomrule
\end{tabular}
\caption{Frozen RH-384 certificate grid.  Finite rows reproduce proved
identities and do not constitute asymptotic evidence.}
\label{tab:certificate}
\end{table}

The mutations include the missing Abel boundary constant, the exponent
$2d$ in place of $2d-\ell$, incorrect partition constants and log powers,
loss of the factor $3$ in the $S_y$ normalizer, the sign $Y+2m$, both
forbidden PNT substitutions in the exact residual, release/hash rebinding,
binary-float injection, Boolean/integer aliasing, duplicate or nonfinite
JSON, and moving the integer cutoff to either neighboring prime.  The two
endpoint mutations are explicitly marked as exact-interface failures only;
they are not counted as failures of the leading PNT equivalent.

All decimal operations run inside fresh precision-$80$ directed contexts.
Three hostile ambient contexts, including altered precision, rounding,
exponent ranges, and trapped underflow/clamping, produce byte-identical
certificates.  The canonical exact certificate has $48689$ bytes and
SHA-256
\begin{center}
\texttt{01c91e57a01de9841f282327ab2f6e1a9368e136393ddab7a2cfe6b019a519c8}.
\end{center}
The closed Draft 2020--12 schema recursively fixes every member, value,
type, and array length.  Strict loaders reject duplicate keys and nonfinite
constants.

The immutable source closure contains $51$ release blobs, with group sizes
$7,8,8,8,8,8,2,2$: RH-374; RH-379 through RH-383; the two-file synthesis
archive; and the RH-2 PNT main text and bibliography.  The aggregate source
digest is
\begin{center}
\texttt{90434e0468ecc062cb522da096a267748725b5dca8e59c642bb7711f45a3e0e4}.
\end{center}
Each live file is checked against its declared release blob before result
generation.  Mutable root policy and handoff files are excluded.  The
publication manifest separately locks the manuscript, PDF, certificate,
schema, audit ledgers, code, tests, and every external input.

\section{Claim boundary and theorem budget}

The standalone verdict is \textbf{Route A: GO}.  The fixed-$r$ Abel
constant, fixed-partition scale compiler, strict scale ordering, five gap
limits, and certified positive intermediate coefficient are new theorem
content inside the frozen class.  The RH-compatibility verdict is
\textbf{Route B: STOP\_SCOPED}.  The following firewalls remain explicit:
\begin{itemize}[leftmargin=2em]
 \item $r$ and $\lambda$ are fixed; no uniform growing-degree theorem or
       effective PNT rate is claimed;
 \item the strict endpoint and successor atom are exact interface data,
       while the leading PNT equivalent alone cannot distinguish their
       endpoint mutations;
 \item all second and lower gap limits retain the exact intrinsic
       subtractions; no bare-PNT surrogate is inserted;
 \item $q$ is fixed before $N\to\infty$; no growing clock $q(N)$, exchange
       of limits, or adaptive-capacity limit is proved;
 \item nonzero phasewise $c_{11}$ remains excluded; no shift-two
       cancellation theorem is inferred;
 \item no intrinsic determinant, scattering completion, self-adjoint
       generator, von Mangoldt weighted prime-power trace, completed-zeta
       divisor equality, or Riemann-zero identification is constructed.
\end{itemize}
Accordingly Gates A--E remain false/open.  This paper neither proves nor
reduces the Riemann hypothesis.  Any successor route must pay for its own
new theorem edge rather than promoting this arithmetic tail translation to
an operator statement.  The four-volume synthesis remains the preserved
foundation \citep{RHMVP2}.

\section*{Declarations}

\noindent\textbf{Data and code availability.}
The analytic proof, exact certificate generator, closed schema, tests,
source-lock ledger, publication manifest, and both byte-identical PDF names
are included with this paper.  The calculation uses no external dataset.

\noindent\textbf{Author contributions.}
The single author is responsible for conceptualization, methodology,
formal analysis, software, validation, writing, and final approval.

\noindent\textbf{Funding.}
No external funding was received for this work.

\noindent\textbf{Competing interests.}
The author declares no competing interests.

\noindent\textbf{Ethics statement.}
No human participants, animals, personal data, or clinical materials were
involved; ethics review and informed consent are not applicable.

\noindent\textbf{AI assistance.}
AI-assisted tools were used for proof auditing, symbolic and numerical code
checking, adversarial test design, and typesetting support.  All claims,
citations, source locks, computations, and final text were checked under the
author's responsibility.

\begingroup
\small
\bibliographystyle{plainnat}
\bibliography{references}
\endgroup

\end{document}
